\section{The Cursor / Claude-Code Plugin}\label{the-cursor-claude-code-plugin}

The harness ships as a single plugin, \textbf{\passthrough{\lstinline!pratyaksha-context-eng-harness!}} (v1.0.0, MIT-licensed), installable into Cursor, Claude Code (CLI and the VS Code extension), and Claude Desktop via the Model Context Protocol \citep{anthropic2025mcp}. This section specifies the shipped artefact in enough detail that an external reviewer can verify its claims directly against the repository at \passthrough{\lstinline!github.com/SharathSPhD/pratyaksha-context-eng-harness!}.

\subsection{Manifest and discovery}\label{manifest-and-discovery}

Two manifests are shipped. \textbf{\passthrough{\lstinline!.claude-plugin/plugin.json!}} is the Claude-Code-format component manifest declaring the 3 agents, 4 commands, 3 skills, and the \passthrough{\lstinline!hooks/hooks.json!} registration. \textbf{\passthrough{\lstinline!marketplace.json!}} is the top-level marketplace manifest used for \passthrough{\lstinline!/plugin install!} discovery, listing the canonical id, version, repository URL, license, keywords, and feature catalogue. Both manifests are version-pinned to \passthrough{\lstinline!1.0.0!}, point to the same upstream URL, and resolve identically in Cursor (via the plugin marketplace \citep{cursor2025plugins}) and Claude Code (via the new plugin marketplace surface \citep{anthropic2025claudecode}). Setup is a single command --- \passthrough{\lstinline!curl -LsSf https://astral.sh/uv/install.sh | sh!} --- and the MCP server's Python dependencies are installed lazily on first tool call by \passthrough{\lstinline!uv!} \citep{uv2025}, so a fresh user can install the plugin without ever invoking \passthrough{\lstinline!pip!}, creating a virtualenv, or running \passthrough{\lstinline!claude mcp add!}.

\subsection{The 15 MCP tools}\label{the-15-mcp-tools}

The MCP server (\passthrough{\lstinline!mcp/server.py!}, FastMCP-style) exposes exactly fifteen tools. Their names, signatures, and operational roles are:

{\def\LTcaptype{none} % do not increment counter
\begin{longtable}[]{@{}
  >{\raggedright\arraybackslash}p{(\linewidth - 6\tabcolsep) * \real{0.2500}}
  >{\raggedright\arraybackslash}p{(\linewidth - 6\tabcolsep) * \real{0.2500}}
  >{\raggedright\arraybackslash}p{(\linewidth - 6\tabcolsep) * \real{0.2500}}
  >{\raggedright\arraybackslash}p{(\linewidth - 6\tabcolsep) * \real{0.2500}}@{}}
\toprule\noalign{}
\begin{minipage}[b]{\linewidth}\raggedright
\#
\end{minipage} & \begin{minipage}[b]{\linewidth}\raggedright
Tool
\end{minipage} & \begin{minipage}[b]{\linewidth}\raggedright
Signature
\end{minipage} & \begin{minipage}[b]{\linewidth}\raggedright
Operationalizes
\end{minipage} \\
\midrule\noalign{}
\endhead
\bottomrule\noalign{}
\endlastfoot
1 & \passthrough{\lstinline!context\_insert!} & \passthrough{\lstinline!InsertInput → ContextElement!} & Avacchedaka insertion (Section 4.2) \\
2 & \passthrough{\lstinline!context\_retrieve!} & \passthrough{\lstinline!RetrieveInput → list[ContextElement]!} & Default avacchedaka query \\
3 & \passthrough{\lstinline!context\_get!} & \passthrough{\lstinline!GetByIdInput → ContextElement!} & Direct id-based fetch (audit) \\
4 & \passthrough{\lstinline!context\_sublate!} & \passthrough{\lstinline!SublateInput → ContextElement!} & Bādha (basic, by id) \\
5 & \passthrough{\lstinline!sublate\_with\_evidence!} & \passthrough{\lstinline!SublateWithEvidenceInput → ...!} & Bādha with explicit pointer + reason \\
6 & \passthrough{\lstinline!detect\_conflict!} & \passthrough{\lstinline!RetrieveInput → \{conflicted: bool, items: ...\}!} & Bayesian-margin conflict check \\
7 & \passthrough{\lstinline!list\_qualificands!} & \passthrough{\lstinline!() → list[str]!} & Diagnostic surface \\
8 & \passthrough{\lstinline!compact!} & \passthrough{\lstinline!CompactInput → CompactionReport!} & Adaptive forgetting (Section 5.7) \\
9 & \passthrough{\lstinline!boundary\_compact!} & \passthrough{\lstinline!BoundaryCompactInput → ...!} & Event-boundary compaction (Section 5.6) \\
10 & \passthrough{\lstinline!context\_window!} & \passthrough{\lstinline!ContextWindowInput → WindowSnapshot!} & Visible-context summary for Manas \\
11 & \passthrough{\lstinline!set\_sakshi!} & \passthrough{\lstinline!SetSakshiInput → ()!} & Pin a session-stable witness invariant \\
12 & \passthrough{\lstinline!get\_sakshi!} & \passthrough{\lstinline!() → list[SakshiInvariant]!} & Retrieve the witness invariants \\
13 & \passthrough{\lstinline!classify\_khyativada!} & \passthrough{\lstinline!ClassifyInput → KhyatiResult!} & 7-class hallucination classifier \\
14 & \passthrough{\lstinline!budget\_status!} & \passthrough{\lstinline!BudgetStatusInput → BudgetStatus!} & TokenBudgetWatchdog (Section 5.8) \\
15 & \passthrough{\lstinline!budget\_record!} & \passthrough{\lstinline!BudgetRecordInput → BudgetStatus!} & Record consumed tokens per turn \\
\end{longtable}
}

Every tool is typed via Pydantic v2 input/output models, returns JSON-serializable structures, and is exercised by at least one unit test in \passthrough{\lstinline!tests/test\_v2/!}. A complete tool reference with full schemas is in \textbf{Appendix B}.

\subsection{The 3 sub-agents}\label{the-3-sub-agents}

\textbf{\passthrough{\lstinline!agents/manas.md!}} defines the \emph{attentional sense-organ} sub-agent: the system prompt instructs the model to call \passthrough{\lstinline!get\_sakshi!}, \passthrough{\lstinline!context\_retrieve!}, \passthrough{\lstinline!context\_window!}, \passthrough{\lstinline!classify\_khyativada!}, and \passthrough{\lstinline!budget\_record!}, and to return structured JSON of the form \passthrough{\lstinline!\{ "draft": "...", "grounding": ["<element\_id>", ...], "uncertain\_claims": ["..."], "needs\_buddhi": true | false \}!} (verbatim contract in the repository file); Manas must \emph{never} emit a user-visible final answer. \textbf{\passthrough{\lstinline!agents/buddhi.md!}} is the \emph{determinative judging} sub-agent, which may call \passthrough{\lstinline!sublate\_with\_evidence!}, \passthrough{\lstinline!detect\_conflict!}, and \passthrough{\lstinline!classify\_khyativada!}, and emits the user-visible answer together with its self-classified \passthrough{\lstinline!khyati\_class!} and \passthrough{\lstinline!confidence!}. \textbf{\passthrough{\lstinline!agents/sakshi-keeper.md!}} is the \emph{witness keeper}: read-only on the live store and append-only on the JSON-lines audit log at \passthrough{\lstinline!\~/.cache/pratyaksha/audit.jsonl!}; it is invoked from the lifecycle hooks (Section 6.5). The choice to ship Manas and Buddhi as \emph{agents} (system-prompt-only sub-agents driven by the host LLM) rather than as fine-tunes is deliberate: it preserves the harness's host-platform-agnosticism (Section 5.4) and lets users swap the underlying model without re-training anything.

\subsection{The 3 skills}\label{the-3-skills}

Following the Cursor/Claude-Code skill format (an \passthrough{\lstinline!SKILL.md!} per directory): \textbf{\passthrough{\lstinline!skills/context-discipline/SKILL.md!}} surfaces the harness's discipline as a triggered skill (\emph{``every retrieved item must be entered with avacchedaka qualifiers; no free-text dumps''}) and activates when the host agent is about to call a search/RAG tool; \textbf{\passthrough{\lstinline!skills/sublate-on-conflict/SKILL.md!}} triggers Buddhi's call to \passthrough{\lstinline!sublate\_with\_evidence!} when \passthrough{\lstinline!detect\_conflict!} returns true; and \textbf{\passthrough{\lstinline!skills/witness-prefix/SKILL.md!}} wraps every Buddhi prompt with a stable \passthrough{\lstinline!<sakshi\_invariants>!} system block via \passthrough{\lstinline!get\_sakshi!}, so the witness frame is \emph{literally a system message} rather than inlined into the user turn.

\subsection{The 4 slash commands}\label{the-4-slash-commands}

The four commands are \textbf{\passthrough{\lstinline!/context-status!}}, which pretty-prints the visible context window by category with a budget gauge; \textbf{\passthrough{\lstinline!/sublate <target\_id> <by\_id> <reason>!}}, the manual sublation override for audit and debug; \textbf{\passthrough{\lstinline!/budget!}}, a shorthand for \passthrough{\lstinline!budget\_status!} rendered as a one-line gauge; and \textbf{\passthrough{\lstinline!/compact-now [strategy]!}}, which manually triggers compaction with \passthrough{\lstinline!strategy ∈ \{adaptive, lru, none\}!} (default \passthrough{\lstinline!adaptive!}).

\subsection{The 3 lifecycle hooks}\label{the-3-lifecycle-hooks}

\passthrough{\lstinline!hooks/hooks.json!} registers three hooks against the standard Claude Code lifecycle events:

\begin{lstlisting}
{
  "hooks": {
    "SessionStart": [{ "matcher": "*", "hooks": [
      { "type": "command", "command": "${CLAUDE_PLUGIN_ROOT}/hooks/session-start.sh", "timeout": 10 }
    ]}],
    "PreToolUse": [{ "matcher": "mcp__pratyaksha_mcp__.*", "hooks": [
      { "type": "command", "command": "${CLAUDE_PLUGIN_ROOT}/hooks/pretooluse-budget.sh", "timeout": 5 }
    ]}],
    "Stop": [{ "matcher": "*", "hooks": [
      { "type": "command", "command": "${CLAUDE_PLUGIN_ROOT}/hooks/stop-compact.sh", "timeout": 5 }
    ]}]
  }
}
\end{lstlisting}

Their roles divide cleanly across the three phases of the turn. \textbf{\passthrough{\lstinline!session-start.sh!}} seeds the session-stable Sākṣī invariants (working directory, git SHA, model id, plugin version) via \passthrough{\lstinline!set\_sakshi!}. \textbf{\passthrough{\lstinline!pretooluse-budget.sh!}} consults \passthrough{\lstinline!budget\_status!} before any harness tool runs and \textbf{logs} over-threshold use by default (advisory hook); if \passthrough{\lstinline!PRATYAKSHA\_BUDGET\_STRICT=1!} is set, the same script can be configured to \textbf{deny} tool execution when the gauge exceeds a hard threshold (default 95\% of \passthrough{\lstinline!context\_budget!}). \textbf{\passthrough{\lstinline!stop-compact.sh!}} invokes \passthrough{\lstinline!compact(strategy="adaptive")!} at the end of every turn so memory pressure does not accumulate across turns.

\subsection{Worked example: a Redis-caching turn, end-to-end}\label{worked-example-a-redis-caching-turn-end-to-end}

To make the runtime contract concrete we trace a single user turn end-to-end through the deployed plugin. The user prompt is \emph{``how do I cache a user session in Redis?''}; the agent's tool returns a mix of pre-Redis-7 blog snippets and the official Redis 7 documentation. Figure\textasciitilde{}\ref{fig:swimlane} visualises the swimlane across \textbf{User → Manas → Buddhi → Sublation → Sākṣī}, with the per-stage immutable JSON line written into \passthrough{\lstinline!\~/.cache/pratyaksha/audit.jsonl!}.

% =====================================================================
% Fig 6 -- End-to-end worked example swimlane (Redis caching case)
% =====================================================================
\begin{figure}[!tbp]
\centering
\resizebox{\linewidth}{!}{%
\begin{tikzpicture}[
  font=\small,
  every node/.style={rounded corners=2pt},
  lane/.style ={draw=black!40,fill=gray!4,inner sep=3mm,minimum height=18mm},
  ev/.style   ={draw=black!60,fill=blue!6,inner sep=3pt,minimum width=30mm,
                minimum height=10mm,align=center,font=\footnotesize},
  badha/.style={draw=black!60,fill=red!10,inner sep=3pt,minimum width=30mm,
                minimum height=10mm,align=center,font=\footnotesize},
  ok/.style   ={draw=black!60,fill=green!10,inner sep=3pt,minimum width=30mm,
                minimum height=10mm,align=center,font=\footnotesize},
  arr/.style  ={-{Latex[length=2mm]},gray!70,thick},
  obs/.style  ={-{Latex[length=2mm]},dashed,violet!70,thick},
  >=Latex
]

% --- Lane labels (left, vertical) ---
\node[font=\footnotesize\bfseries,anchor=east] at (-2,  0)   {User};
\node[font=\footnotesize\bfseries,anchor=east] at (-2, -2.0) {Manas};
\node[font=\footnotesize\bfseries,anchor=east] at (-2, -4.0) {Buddhi};
\node[font=\footnotesize\bfseries,anchor=east] at (-2, -6.0) {Sublation};
\node[font=\footnotesize\bfseries,anchor=east] at (-2, -8.0) {Sak\d{s}\=\i};

% --- Lane backgrounds ---
\foreach \y/\h in {0/1.4, -2.0/1.4, -4.0/1.4, -6.0/1.4, -8.0/1.4}{
  \draw[draw=black!20,fill=gray!3] (-1.7,\y-\h/2) rectangle (16.5,\y+\h/2);
}

% --- Events (left to right) ---
\node[ev]    (e1) at ( 0.5,  0)    {$t_1$: ``cache user\\session in Redis''};
\node[ev]    (e2) at ( 4.5, -2.0)  {$t_2$: retrieve\\3 stale items};
\node[badha] (e3) at ( 8.5, -4.0)  {$t_3$: TYPE\_CLASH\\(blog vs.\ docs)};
\node[ok]    (e4) at (12.5, -6.0)  {$t_4$: docs win,\\blog sublated};
\node[ok]    (e5) at (15.0, -4.0)  {$t_5$: tagged answer\\(yath\=artha)};

% --- Sakshi log entries (bottom lane) ---
\node[ev,fill=violet!4,font=\scriptsize\ttfamily,minimum width=24mm]
   (l1) at ( 0.5, -8.0) {evt=USER\_TURN};
\node[ev,fill=violet!4,font=\scriptsize\ttfamily,minimum width=24mm]
   (l2) at ( 4.5, -8.0) {evt=ATTEND\\K=3 budget=1.2k};
\node[ev,fill=violet!4,font=\scriptsize\ttfamily,minimum width=24mm]
   (l3) at ( 8.5, -8.0) {evt=BADHA\_DETECT\\kind=TYPE\_CLASH};
\node[ev,fill=violet!4,font=\scriptsize\ttfamily,minimum width=24mm]
   (l4) at (12.5, -8.0) {evt=BADHA\_RESOLVE\\sublated=item\_71};
\node[ev,fill=violet!4,font=\scriptsize\ttfamily,minimum width=24mm]
   (l5) at (15.0, -8.0) {evt=ANSWER\\class=yath\=artha};

% --- Pipeline arrows ---
\draw[arr] (e1) -- (e2);
\draw[arr] (e2) -- (e3);
\draw[arr] (e3) -- (e4);
\draw[arr] (e4) -- (e5);

% --- Sakshi observation arrows (passive, dashed) ---
\foreach \src/\dst in {e1/l1, e2/l2, e3/l3, e4/l4, e5/l5}
  \draw[obs] (\src) -- (\dst);

\end{tikzpicture}}
\caption{\textbf{End-to-end worked example: a Redis-caching turn.}
At $t_1$ the user asks how to cache a session in Redis. \textbf{Manas}
($t_2$) retrieves three stale items from the context store under a
1.2k-token budget. \textbf{Buddhi} ($t_3$) detects a TYPE\_CLASH between
a blog post and the official Redis 7 documentation sharing the same
$\langle$Redis-session, expiry-policy$\rangle$ qualificand--qualifier
slot. The \textbf{sublation} step ($t_4$) ranks sources by limitor
precedence: the official docs win, the blog item is marked sublated.
At $t_5$ Buddhi emits the tagged answer with class $5$ (yath\=artha,
veridical). The \textbf{Sak\d{s}\=\i Keeper} (bottom lane) writes
exactly one immutable JSON line per stage; nothing on the response path
is touched, but the entire turn is independently replayable.}
\label{fig:swimlane}
\end{figure}


The five host-visible artefacts are: (i) one \passthrough{\lstinline!mcp\_\_pratyaksha\_mcp\_\_manas\_step!} JSON-RPC call (Manas attends, \(K=3\) items under a 1.2 K-token budget); (ii) one \passthrough{\lstinline!detect\_conflict!} round-trip flagging a TYPE\_CLASH on the \passthrough{\lstinline!(Redis-session, expiry-policy)!} qualificand-qualifier slot; (iii) one \passthrough{\lstinline!sublate\_with\_evidence!} call that retires the blog post in favour of the Redis 7 docs (limitor precedence: official docs \passthrough{\lstinline!prec=8!} \textgreater{} blog post \passthrough{\lstinline!prec=2!}); (iv) one \passthrough{\lstinline!mcp\_\_pratyaksha\_mcp\_\_buddhi\_step!} call returning the answer with \passthrough{\lstinline!khyati\_class = "yathārtha"!} (veridical) and \passthrough{\lstinline!confidence = 0.91!}; and (v) one append to the Sākṣī log per stage. The complete \passthrough{\lstinline!/context-status!} snapshot, the raw JSON-RPC payloads, and the four matching audit-log lines are reproduced verbatim in \textbf{Appendix C.8} so that the turn can be replayed byte-for-byte from the shipped plugin against the cached evidence trail.

This single turn exercises every load-bearing primitive of Sections 4--5 in a non-coding agentic context, and is the runtime template the L1 (§8) and L3 (§10) experiments instantiate at scale.

\subsection{\texorpdfstring{What the plugin does \emph{not} contain}{What the plugin does not contain}}\label{what-the-plugin-does-not-contain}

We assert and audit two negative claims. \textbf{First, zero references to the dev-time orchestration tools} \passthrough{\lstinline!attractor-flow!}, \passthrough{\lstinline!ralph-loop!}, and \passthrough{\lstinline!triz-engine!}: those tools were used to build the harness (Section 3) but are not runtime dependencies, a grep over the entire shipped plugin tree returns zero matches, and the audit is preserved as \passthrough{\lstinline!docs/no\_dev\_deps\_audit.md!} in the repository. \textbf{Second, zero hard dependencies on heavyweight stacks} like \passthrough{\lstinline!vllm!}, \passthrough{\lstinline!mlflow!}, \passthrough{\lstinline!chromadb!}, or \passthrough{\lstinline!huggingface-hub!}: the harness's runtime requirements are \passthrough{\lstinline!mcp!}, \passthrough{\lstinline!pydantic!}, \passthrough{\lstinline!tiktoken!}, \passthrough{\lstinline!numpy!}, and \passthrough{\lstinline!anthropic!}, and the optional Qwen3-1.7B surprise backend is loaded only when the user opts in. This matters because the harness is a \emph{discipline}, not an infrastructure. Anything heavier than that would defeat the purpose of a plugin you can install in 30 seconds.

\subsection{Smoke test and CI}\label{smoke-test-and-ci}

\passthrough{\lstinline!mcp/smoke\_test.py!} exercises every tool against the local MCP server: 15 round-trips, each verifying the input schema, the output schema, and one non-trivial behavioural assertion (e.g.~\passthrough{\lstinline!sublate\_with\_evidence!} actually flips the target item's status). The smoke test runs in roughly 4 seconds wall-clock and is wired into the repo's CI as \passthrough{\lstinline!pytest -m smoke!}. We additionally exercise every command and hook via \passthrough{\lstinline!claude --debug!} in the developer's local Claude Code installation; the transcript of one such smoke run is preserved in \passthrough{\lstinline!docs/plugin\_smoke\_transcript.md!}.

\subsection{Hot-swappability across hosts}\label{hot-swappability-across-hosts}

The same \passthrough{\lstinline!marketplace.json!} resolves in \textbf{Cursor} (via the IDE's plugin browser, registered by URL), in \textbf{Claude Code (CLI)} (via \passthrough{\lstinline!/plugin install <url>!}), in \textbf{Claude Code (VS Code extension)} (via the same \passthrough{\lstinline!/plugin install!} UI), and in \textbf{Claude Desktop} (via the Desktop app's MCP-based plugin surface). We have manually exercised the install path on all four hosts during development. The harness's runtime invariants (Sections 4 and 5) hold identically across all four because the only inter-process surface is MCP, and MCP is host-agnostic by construction \citep{anthropic2025mcp}.
